\documentclass[11pt]{article}
\usepackage[a4paper,margin=2cm]{geometry}
\usepackage{amsmath,amssymb}
\usepackage[hidelinks]{hyperref}
\allowdisplaybreaks
\title{Qpu.Processing.Verify.Hexbit}
\author{Tsvetan Rouschev}
\date{@uuidna/qpu v1.1.1, CC-BY-NC-ND-4.0}
\begin{document}
\maketitle
\noindent\url{https://github.com/uuidna/qpu/blob/v1.1.1/lean/Qpu/Processing/Verify/Hexbit.lean}

\paragraph{\texttt{Qpu.Verify.Hexbit.gap\_is\_rejected}}
Control: the ASCII code of ':' (58) sits in the gap between digits and letters and is likewise rejected.
\[ \mathrm{Qpu.Verify.Hexbit.unhexCode}\ 58 = 256 \]
\noindent axioms: none; proof: \url{https://github.com/uuidna/qpu/blob/v1.1.1/lean/Qpu/Processing/Verify/Hexbit.lean\#L50}

\paragraph{\texttt{Qpu.Verify.Hexbit.hex\_roundtrips}}
Every nibble survives the round trip: encode then decode is the identity on 0..15, enumerated case by case so the kernel decides each limb.
\[ \forall n < 16, \mathrm{Qpu.Verify.Hexbit.unhexCode} (\mathrm{Qpu.Verify.Hexbit.hexCode}\ n) = n \]
\noindent axioms: \texttt{propext}, \texttt{Quot.sound}; proof: \url{https://github.com/uuidna/qpu/blob/v1.1.1/lean/Qpu/Processing/Verify/Hexbit.lean\#L25}

\paragraph{\texttt{Qpu.Verify.Hexbit.non\_hex\_is\_rejected}}
Control: the ASCII code of 'g' (103) does not decode to a nibble.
\[ \mathrm{Qpu.Verify.Hexbit.unhexCode}\ 103 = 256 \]
\noindent axioms: none; proof: \url{https://github.com/uuidna/qpu/blob/v1.1.1/lean/Qpu/Processing/Verify/Hexbit.lean\#L46}

\paragraph{\texttt{Qpu.Verify.Hexbit.prefix\_is\_high\_half}}
Routing: the 32-bit hexbit prefix of a 64-bit id is its high half.
\[ 1311768467294899695 / 4294967296 = 305419896 \]
\noindent axioms: none; proof: \url{https://github.com/uuidna/qpu/blob/v1.1.1/lean/Qpu/Processing/Verify/Hexbit.lean\#L53}

\paragraph{\texttt{Qpu.Verify.Hexbit.prefix\_width}}
Control: the prefix really is 8 hexits — 32 bits — wide.
\[ 305419896 < 2 ^{(8 \cdot 4)} \]
\noindent axioms: none; proof: \url{https://github.com/uuidna/qpu/blob/v1.1.1/lean/Qpu/Processing/Verify/Hexbit.lean\#L57}
\end{document}
